% % Reemplace por el título de su capítulo.
\chapter{El regreso}

% % Reemplace por un epígrafe genérico y su atribución.
\epigraph{Volver a casa es empezar otra vez.}{Atribuido a una copla anónima}

% % Reemplace por la escena inicial: el camino de vuelta.
El regreso empieza con la maleta ligera y el billete en el bolsillo. Por
la ventana pasan campos cortos y casas bajas que parecen contar los
minutos. Nadie espera nada grande: solo llegar, saludar y sentarse a la
mesa. Y en ese gesto simple cabe la frase que sostiene la casa,
\enquote{lo bueno tarda, pero llega}.

% % Reemplace por el desarrollo: qué encuentra al volver.
Al llegar, todo está en su lugar aunque algo ha cambiado de sitio. La
cocina huele a pan, el patio guarda la misma silla y el reloj marca la
hora de siempre. Las noticias se cuentan despacio, sin orden, mientras la
tarde cae sobre los tejados.

\begin{center}
  * * *
\end{center}

% % Reemplace por el cierre: qué se queda después del regreso.
Por la noche no hay discurso ni promesa, solo un acuerdo claro: cuidar lo
cercano y volver a empezar cada día. El libro se cierra aquí, con la casa
en silencio y la lámpara encendida para quien llegue después.
